\documentclass[a4paper,10pt]{article}

\usepackage[utf8]{inputenc}
\usepackage[T1]{fontenc}
\usepackage[ngerman]{babel}
\usepackage{amsmath}
\usepackage{amssymb}
\usepackage{xcolor}
\usepackage{enumitem}
\usepackage{geometry}
\geometry{a4paper,margin=1.9cm}

% ----- White Mode -----
\definecolor{accent}{HTML}{1A4E8A}   % dunkelblau (Kapitel)
\definecolor{accent2}{HTML}{8A5A00}  % gold/braun (Aufgaben)
\definecolor{accent3}{HTML}{2E6E3A}  % grün (Selbstreflexion)
\definecolor{accent4}{HTML}{7A3E9A}  % violett (Denkanstoß)
\definecolor{rulecol}{HTML}{B0B0B0}

\newcommand{\kapitel}[1]{%
  \par\vspace{7pt}
  {\large\bfseries\color{accent} #1}\par
  \noindent\textcolor{accent}{\rule{\linewidth}{0.8pt}}\par\vspace{2pt}}

\newcommand{\aufg}[2]{%
  \par\vspace{3pt}
  {\bfseries\color{accent2} #1}~~{\footnotesize\color{accent2}[#2]}\par\vspace{1pt}}

\newcommand{\teaser}[1]{%
  \par\vspace{2pt}{\color{accent4}\textbf{Denkanstoß.}}~#1\par}

\newcommand{\marg}[1]{\par\vspace{1pt}{\footnotesize\itshape Randnotiz: #1}\par}

\newcommand{\reflexion}{%
  \par\vspace{2pt}{\bfseries\color{accent3} Selbstreflexionsfragen}\par\vspace{1pt}}

\setlist[enumerate]{leftmargin=1.8em,itemsep=1pt,topsep=1pt,parsep=0pt}
\setlist[itemize]{leftmargin=1.3em,itemsep=1pt,topsep=1pt,parsep=0pt,label=\textcolor{accent3}{\textbullet}}
\setlength{\parindent}{0pt}
\setlength{\parskip}{3pt}

\begin{document}

\begin{center}
  {\Large\bfseries\color{accent} Numerische Methoden -- Übungsaufgaben}\\[1pt]
  {\small Vollständig aus dem Vorlesungsskript \emph{``Numerical Methods''}
   (Prof.\ Dr.-Ing.\ Mark Schutera) extrahiert \;·\; ins Deutsche übersetzt}
\end{center}
\noindent\textcolor{rulecol}{\rule{\linewidth}{0.4pt}}

{\footnotesize Diese Sammlung enthält \emph{alle} im Skript vorkommenden Aufgaben
und durchgerechneten Beispiele -- Handrechnungen, Code-/Maschinen-Aufgaben,
Margin-Aufgaben, Denkanstöße und Selbstreflexionsfragen -- über alle sieben
Kapitel. Lösungswege: siehe \texttt{loesungen}.}

% =====================================================================
\kapitel{Kapitel 1 \;--\; Einstieg in die Numerischen Methoden}

\aufg{1.1 \;--\; Minimum der Sinusfunktion}{Handrechnung}
Bestimmen Sie das Minimum der Funktion $\sin(x)$ auf dem Intervall $[-3,1]$, indem
Sie die Funktion mit Schrittweite $h$ diskretisieren, die Funktionswerte an den
Gitterpunkten auswerten und den kleinsten auswählen. Vergleichen Sie mit dem
analytisch exakten Minimum.

\aufg{1.2 \;--\; Grid-Search}{Handrechnung}
Führen Sie eine Grid-Search für $w$ und $b$ mit Schrittweite $2{,}5$ über den
Bereich $w,b\in\{0;\,2{,}5;\,5;\,7{,}5;\,10\}$ durch. Das zugrunde liegende System
ist $\big(\begin{smallmatrix}2&1\\5&1\end{smallmatrix}\big)\big(\begin{smallmatrix}w\\b\end{smallmatrix}\big)=\big(\begin{smallmatrix}11\\13\end{smallmatrix}\big)$.
Berechnen Sie für jede Kombination den Fehler $\mathrm{error}_1=|2w+b-11|$,
$\mathrm{error}_2=|5w+b-13|$. Werten Sie alle Kombinationen aus und wählen Sie das
Paar $(w,b)$ mit dem kleinsten akkumulierten Fehler.

\aufg{1.3 \;--\; Enter the machine}{Code}
Viele numerische Verfahren sind für die Umsetzung am Computer gedacht. Sie erhalten
ein vorgehostetes Jupyter-Notebook mit einer Python-Vorlage für diese Übung. Nutzen
Sie die Selbstreflexionsfragen, um sich beim Erkunden und Experimentieren mit der
Implementierung leiten zu lassen.

\teaser{Fällt Ihnen eine einfache Möglichkeit ein, das Verhältnis von Genauigkeit
zu Rechenaufwand unseres Grid-Search-Beispiels zu verbessern?}

\reflexion
\begin{itemize}
  \item Warum ist die Wahl der Schrittweite $h$ beim Diskretisieren einer stetigen
        Funktion wichtig, und wie beeinflusst sie Genauigkeit und Rechenzeit?
  \item Wie beeinflusst die Wahl des Intervalls $[a,b]$ die Ergebnisse der
        Diskretisierung und die Lage der numerisch gefundenen Extrema?
  \item Was sind die wesentlichen Unterschiede zwischen einer stetigen Funktion und
        ihrer diskretisierten Version, und welche Implikationen ergeben sich daraus?
\end{itemize}

% =====================================================================
\kapitel{Kapitel 2 \;--\; Gleitkomma-Arithmetik}

\aufg{2.1 \;--\; Rundungsfehler}{Handrechnung}
Bekommen wir an einem einfachen Beispiel ein Gefühl für den Rundungsfehler.
Betrachten Sie die Division von $10$ durch $3$. Schreiben Sie die Dezimalentwicklung
aus und bestimmen Sie den Rundungsfehler beim Abschneiden nach drei
Nachkommastellen.

\aufg{2.2 \;--\; 2-Bit-Mantisse \& Maschinenepsilon}{Handrechnung}
Eine unnormalisierte Mantisse hat $2$ Bit, der Signifikand ist $0{.}xx_2$. Bestimmen
Sie alle darstellbaren Werte und das zugehörige Maschinenepsilon
$\varepsilon_{\mathrm{mach}}=2^{-t}$. Was bedeutet das für relative bzw.\ absolute
Präzision?

\aufg{2.3 \;--\; Auslöschung von Hand}{Handrechnung}
Zwei Zahlen werden mit $8$ signifikanten Stellen gespeichert.
\textbf{(a)} $a=12345678{,}5$, $b=12345678{,}0$: bestimmen Sie $\tilde a,\tilde b$
und $\tilde a-\tilde b$ und vergleichen Sie mit $a-b$.
\textbf{(b)} $a=12345678{,}4$, $b=12345678{,}0$: ebenso. Was beobachten Sie beim
relativen Fehler?

\aufg{2.4 \;--\; Festkomma-Skalierung}{Handrechnung}
Sie speichern Zahlen aus $[-1000,1000]$ in einem 32-Bit-Integer mit Skalierungsfaktor
$10^{6}$. Wie wird $1{,}234567$ gespeichert? Welche Präzision (kleinste darstellbare
Differenz) und welcher maximale Rundungsfehler ergeben sich?

\aufg{2.5 \;--\; Maschinenepsilon (hands on)}{Code}
Überlegen Sie, wie Sie das Maschinenepsilon eines beliebigen Systems für ein
bestimmtes Gleitkommaformat per Programm bestimmen würden. Rufen Sie sich die
Definition in Erinnerung und schreiben Sie Ihre Überlegungen als Pseudocode auf. Was
würde ein solches Programm auf einem Festkomma- gegenüber einem Gleitkommasystem
zeigen? Probieren Sie verschiedene Datentypen am Rechner aus. Wenn Sie einen
Taschenrechner bauen würden -- welchen Datentyp würden Sie wählen und warum?

\aufg{2.6 \;--\; Auslöschung demonstrieren (hands on)}{Code}
Überlegen Sie, wie Sie Auslöschung (loss of significance) an Ihrem Rechner
demonstrieren würden. Schreiben Sie Ihre Überlegungen als Pseudocode auf, bevor Sie
weiterlesen.

\aufg{2.7 \;--\; Argumentation}{Konzept}
Überlegen Sie, wie das Berechnen von $f(\epsilon)=\dfrac{1}{a+\epsilon-b}$ für
kleines $\epsilon$ und $a=b$ dies leisten würde. Was erwarten Sie zu sehen, wenn
$\epsilon$ sehr klein ist?
\marg{Was passiert für $a>b$? Denken Sie in signifikanten Stellen.}

\teaser{Fallen Ihnen Metriken für numerische Verfahren ein, basierend auf den
Approximationen und Fehlern, die wir besprochen haben?}

\reflexion
\begin{itemize}
  \item Wie ist die Gleitkommadarstellung aufgebaut, und was sind ihre Komponenten?
  \item Was ist das Maschinenepsilon, und wie hängt es mit der numerischen
        Genauigkeit zusammen?
  \item Wie wirken sich diese Konzepte in der Praxis auf numerische Berechnungen aus?
\end{itemize}

% =====================================================================
\kapitel{Kapitel 3 \;--\; Fehleranalyse}

\aufg{3.1 \;--\; Minimum der Sinusfunktion (Eigenschaften)}{Handrechnung}
Bestimmen Sie das Minimum von $\sin(x)$ auf $[-3,1]$ noch einmal. Diskretisieren Sie
mit zwei Schrittweiten $h=1$ und $h=0{,}2$. Nehmen Sie an, dass die Eingabedaten $x$
durch Messfehler um $\tilde x=x\pm0{,}25$ gestört sind (zwei Dezimalstellen genau).
Analysieren Sie \textbf{Kondition}, \textbf{Stabilität}, \textbf{Konsistenz} und
\textbf{Konvergenz} des numerischen Verfahrens in beiden Fällen und quantifizieren Sie
sie mit den im vorigen Abschnitt angegebenen Formeln.

\aufg{3.2 \;--\; Stabilität und Kondition}{Beweis}
Zeigen Sie: Die Stabilität entspricht der Kondition des Systems für $h\to0$.

\aufg{3.3 \;--\; Compute-getriebene Fehleranalyse (hands on)}{Code}
Brute-forcen wir uns durch die Fehleranalyse beim Lösen des Zwei-Gleichungs-Systems
aus Kapitel 1. Bestimmen und optimieren Sie Kondition, Stabilität, Konsistenz und
Konvergenz der numerischen Lösung, indem Sie das numerische Verfahren optimieren.

\reflexion
\begin{itemize}
  \item Wie gut konditioniert ist das Zwei-Gleichungs-System verglichen mit der
        numerischen Lösung des Minimums der Sinusfunktion?
  \item Konvergiert die Stabilität beim Zwei-Gleichungs-System mit kleineren
        Schrittweiten? Wogegen?
  \item Wie wirken sich Störungen der Eingabedaten auf die Lösung aus? Gibt es ein
        Zusammenspiel mit der Schrittweite?
  \item Vergleichen Sie Konvergenz mit Stabilität und Konsistenz. Können Sie Ihre
        Beobachtungen in Begriffen von Bias und Varianz ausdrücken?
  \item Beobachten Sie für immer kleinere Schrittweiten eine Grenze der Genauigkeit?
        Wenn ja, warum?
  \item Welches weitere Problem beobachten Sie beim Verfeinern der Schrittweite?
\end{itemize}

\teaser{Bisher haben wir Brute-Force-Lösungen verwendet. Fällt Ihnen eine
Möglichkeit ein, Brute-Force-Methoden zu verfeinern, um mit weniger Aufwand bessere
Ergebnisse zu erzielen?}

% =====================================================================
\kapitel{Kapitel 4 \;--\; Newton-Verfahren}

\aufg{4.1 \;--\; Newton für $\sqrt{2}$}{Handrechnung}
Berechnen Sie $\sqrt{2}$ mit dem Newton-Verfahren, ausgehend von $x_0=2$ (also die
Nullstelle von $f(x)=x^2-2$). Führen Sie einige Iterationen durch und beobachten Sie,
wie sich der Fehler pro Schritt etwa quadriert.

\aufg{4.2 \;--\; Newton vs.\ Grid-Search für $\sin(x)=0$}{Handrechnung}
Vergleichen Sie Grid-Search und Newton-Verfahren beim Finden der Nullstelle von
$\sin(x)=0$ auf $[2{,}5;\,4]$ (Lösung nahe $\pi$). \textbf{(a)} Grid-Search mit
Schrittweite $h=0{,}3$. \textbf{(b)} Newton ($f'=\cos$), Start $x_0=3{,}5$.
\textbf{(c)} Vergleichen Sie Genauigkeit und Anzahl der Auswertungen.

\aufg{4.3 \;--\; Denksportaufgabe: Versagensmodi}{Schaubild}
Das Newton-Verfahren kann auf verschiedene Weisen versagen. Visualisieren Sie die
Versagensmodi im Schaubild von $f(x)=x^3-x$.

\aufg{4.4 \;--\; Newton von Hand}{Handrechnung}
Finden Sie die Nullstelle von $f(x)=x^3-x=0$, indem Sie das Newton-Verfahren
geometrisch und anschließend von Hand anwenden. Die Nullstellen liegen bei
$x=-1,0,1$. Suchen wir die Nullstelle bei $x=1$, ausgehend von $x_0=1{,}4$.

\aufg{4.5 \;--\; Sekantenverfahren von Hand}{Handrechnung}
Wenden Sie das Sekantenverfahren auf dasselbe Problem mit $x_0=1{,}2$, $x_1=1{,}4$
an. Wieder zuerst geometrisch, dann von Hand.

\aufg{4.6 \;--\; Geometrische Fehlersuche}{Handrechnung}
Finden Sie Startpunkte $x_0$ für verschiedene Versagensmodi: einen, bei dem Newton
wegen verschwindender Ableitung versagt, einen, bei dem es zu einer Nicht-Ziel-Nullstelle
konvergiert, und einen, bei dem es zwischen Punkten oszilliert.

\aufg{4.7 \;--\; Enter the machine}{Code}
Betrachten wir Konvergenz und Konvergenzraten von Grid-Search, Newton- und
Sekantenverfahren in Python. Experimentieren Sie mit verschiedenen Startpunkten und
beobachten Sie die Versagensmodi -- zuerst geometrisch, dann analytisch, dann im Code.

\reflexion
\begin{itemize}
  \item Was ist der Hauptvorteil des Newton-Verfahrens gegenüber der Grid-Search?
  \item Warum genügt eine lineare Approximation zum Finden von Nullstellen? Was sagt
        die Taylor-Entwicklung darüber?
  \item Warum ist quadratische Konvergenz so mächtig? Wie viele Iterationen bräuchte
        Newton, um von Fehler $10^{-1}$ auf $10^{-16}$ zu kommen?
  \item In welchen Situationen würden Sie das Sekantenverfahren vorziehen? Warum nicht
        in anderen?
\end{itemize}

\teaser{Wie können wir sicherstellen, dass wir die globale Lösung finden und nicht
nur eine lokale?}

% =====================================================================
\kapitel{Kapitel 5 \;--\; Globale Optimierung}

\aufg{5.1 \;--\; Newton auf Shekel, weiterer Startpunkt}{Code}
Wenden Sie das Newton-Verfahren auf die 1D-Shekel-Funktion von verschiedenen
Startpunkten an und beobachten Sie, wohin es konvergiert. Probieren Sie selbst einen
weiteren Startpunkt und diskutieren Sie das allgemeine Verhalten.
\marg{Wie würden Sie die Update-Regel abändern, um stattdessen Maxima zu finden?}

\aufg{5.2 \;--\; Random Restarts}{Handrechnung}
Angenommen, Sie optimieren eine Funktion mit $5$ lokalen Minima, und das Einzugsgebiet
des globalen Minimums deckt $15\%$ des Suchraums ab ($p=0{,}15$). Wie groß ist die
Wahrscheinlichkeit, das globale Minimum mit $K=10$ Random Restarts zu finden?
Nutzen Sie $P=1-(1-p)^K$ und lösen Sie nach $K$ auf:
$K=\frac{\ln(1-P)}{\ln(1-p)}$. Wie viele Restarts brauchen wir für $99\%$? Berechnen
Sie dann $K$ für $p=0{,}05$ und $p=0{,}01$ beim selben $99\%$-Ziel.

\aufg{5.3 \;--\; Eigene 1D-Shekel-Funktion}{Code}
Entwerfen Sie an Ihrem Rechner Ihre eigene 1D-Shekel-Funktion (anfangs einfach, später
komplexer). Wenden Sie die lokalen Optimierungsverfahren an und zeigen Sie, wo Newton
(erneut) versagt. Untersuchen Sie verschiedene $x_0$ und wie sich das
Konvergenzverhalten ändert.

\aufg{5.4 \;--\; Die Suche zu Minima lenken}{Konzept + Code}
Erinnern Sie sich an den $|f''|$-Trick, der das Vorzeichen der Krümmung im Nenner
umkehrt und den Schritt zwingt, stets bergab zu gehen. Überlegen Sie, was das
geometrisch für die Tangentenkonstruktion bedeutet, und probieren Sie es im Code.

\aufg{5.5 \;--\; Einzugsgebiet (Basin of Attraction)}{Schaubild}
Betrachten Sie die Plots der Shekel-Funktion, markieren Sie die Einzugsgebiete jedes
lokalen Minimums und schätzen Sie ihre relativen Größen. Welches ist das globale
Minimum? Wie hängt das mit der Wahrscheinlichkeit zusammen, es per Random Restarts zu
finden?

\aufg{5.6 \;--\; Lokalen Minima entkommen}{Code}
Setzen Sie Random Restarts und Basin Hopping auf Ihrer 1D-Shekel-Funktion ein und
vergleichen Sie, wie viele Funktionsiterationen jede Methode bis zum globalen Minimum
braucht. Wie anfällig sind sie fürs Hängenbleiben? Wie beeinflusst die
Schrittweitenverteilung die Performance des Basin Hopping?

\aufg{5.7 \;--\; Nicht-konvexe Becken}{Konzept}
Shekel ist im Allgemeinen nicht-konvex, die einzelnen Becken jedoch schon; notieren
Sie eine Funktion, bei der die Einzugsgebiete nicht so leicht abzugrenzen sind.

\aufg{5.8 \;--\; Adaptive Schrittweite}{Konzept + Code}
Fällt Ihnen ein geschickter Weg ein, die Schrittweitenverteilung während der
Optimierung automatisch anzupassen? Was wäre eine gute Heuristik? Wie könnten Sie
katastrophale Sprünge verhindern? Implementieren Sie es.

\aufg{5.9 \;--\; Noise and landscape engineering}{Konzept}
Die Loss-Landschaft ist nicht statisch, sondern soll geformt werden. Überlegen Sie,
wie das Approximieren der Loss-Landschaft auf Mini-Batches (unterschiedliche Auswahl
von Punkten) wirkt.

\aufg{5.10 \;--\; Code-Übung: Restarts vs.\ Basin Hopping}{Code}
Implementieren Sie Random Restarts und Basin Hopping auf der 1D-Shekel-Funktion.
Vergleichen Sie, wie viele Funktionsauswertungen jede Methode braucht, um das globale
Minimum bei $x\approx4$ zu finden.

\aufg{5.11 \;--\; Einzugsgebiete kartieren}{Handrechnung}
Betrachten Sie die 1D-Shekel-Funktion; ein lokaler Optimierer konvergiere stets zum
nächstgelegenen Foxhole.
\textbf{(a)} Skizzieren Sie (grob) die Einzugsgebiete für die drei Foxholes bei
$x=0,4,7$. Wo liegen die Basin-Grenzen?
\textbf{(b)} Wenn Basin Hopping einen Sprung $\delta\sim\mathrm{Uniform}(-3,3)$
ausgehend vom lokalen Minimum $x^{\circ}=7$ verwendet -- wie groß ist die
Wahrscheinlichkeit, dass ein einzelner Sprung im Becken des globalen Minimums landet?
\textbf{(c)} Warum könnte Basin Hopping das globale Minimum hier schneller finden als
Random Restarts?

\reflexion
\begin{itemize}
  \item Warum versagt lokale Optimierung auf nicht-konvexen Landschaften wie Shekels
        Foxholes?
  \item Wie skaliert der erforderliche Aufwand mit der Größe des globalen Beckens bei
        Random Restarts?
  \item Wie unterscheidet sich Basin Hopping von Random Restarts?
  \item Was sagt uns $f''(x_n)$ über die Loss-Landschaft, und wie können wir es
        nutzbar machen?
  \item Wie hilft das Approximieren der Loss-Landschaft mit verrauschten Schätzungen
        (Mini-Batches), lokalen Optima zu entkommen?
  \item Wählen Sie $x_0$ im Becken eines lokalen Minimums; finden Sie Wege, mit denen
        Ihr Verfahren entkommen kann -- wie wirksam ist welcher?
\end{itemize}

\teaser{Was, wenn wir einer hochfrequenten Version von Shekels Foxholes
gegenüberstehen? Welches Problem tritt dabei auf?}

% =====================================================================
\kapitel{Kapitel 6 \;--\; Numerische Integration}

\aufg{6.1 \;--\; Exakter Wert}{Handrechnung}
Berechnen Sie das Integral von $\sin(x)$ von $-2$ bis $-1$ exakt -- als Referenzwert
zur Fehlerquantifizierung der folgenden Methoden.

\aufg{6.2 \;--\; Mittelpunktsregel von Hand}{Handrechnung}
Berechnen Sie $\int_{-2}^{-1}\sin(x)\,dx$ mit der Mittelpunktsregel mit $n=1$
Intervall.

\aufg{6.3 \;--\; Mittelpunktsregel mit mehr Intervallen}{Handrechnung}
Berechnen Sie die Mittelpunktsregel mit $n=2$ Intervallen (Mittelpunkte bei $-1{,}75$
und $-1{,}25$).
\marg{Berechnen Sie die Approximation auch an den Endpunkten des Intervalls und
vergleichen Sie den Fehler.}

\aufg{6.4 \;--\; Trapezregel von Hand}{Handrechnung}
Berechnen Sie mit der Trapezregel $\int_{-2}^{-1}\sin(x)\,dx$ mit $n=2$ Intervallen.
Die Stützstellen sind $-2,\,-1{,}5,\,-1$.

\aufg{6.5 \;--\; Simpson-Regel von Hand}{Handrechnung}
Berechnen Sie mit der Simpson-Regel $\int_{-2}^{-1}\sin(x)\,dx$ mit $n=2$ Intervallen.
Die Stützstellen sind $-2,\,-1{,}5,\,-1$.
\marg{Probieren Sie mehr Intervalle und prüfen Sie sich mit zusammengesetzten
Simpson- und Trapezregeln.}

\aufg{6.6 \;--\; Monte Carlo von Hand}{Handrechnung}
Schätzen Sie $\int_{-2}^{-1}\sin(x)\,dx$ per Monte Carlo mit $N=4$ Zufallsstichproben
(etwa der Rechenaufwand der Trapezmethode). Die Stichproben sind $X_1=-1{,}95$,
$X_2=-1{,}55$, $X_3=-1{,}15$, $X_4=-1{,}05$.

\aufg{6.7 \;--\; Monte Carlo: Fehler \& $N$ erhöhen}{Konzept}
Bestimmen Sie Stichprobenvarianz und Standardfehler des Monte-Carlo-Schätzers. Was
passiert, wenn Sie $N$ auf $16$ erhöhen? Erläutern Sie kurz, warum der Fehler des
Monte-Carlo-Schätzers dimensionsunabhängig ist.

\aufg{6.8 \;--\; Simpson-Gewichte via Lagrange}{Herleitung}
Leiten Sie die Quadraturgewichte der Simpson-Regel her, indem Sie die
Lagrange-Basispolynome über $[a,b]$ integrieren.
\marg{Beginnen Sie mit $x=0$ und sehen Sie, wie die Linearfaktoren ausfallen. Was
passiert, wenn Sie eine lineare Funktion mit einer quadratischen approximieren?}

\aufg{6.9 \;--\; Höhere Dimensionen (on your machine)}{Code}
Implementieren Sie Monte-Carlo-Integration für ein $d$-dimensionales Integral (z.\,B.\
eine multivariate Gauß-Funktion über einen Hyperwürfel). Erhöhen Sie $d$ langsam und
beobachten Sie das Fehlerverhalten und wie gitterbasierte Quadratur rechnerisch
undurchführbar wird.

\reflexion
\begin{itemize}
  \item Wie vergleichen sich die Fehler von Mittelpunkts-, Trapez-, Simpson-Regel und
        Monte Carlo?
  \item Warum hat die Mittelpunktsregel eine bessere Genauigkeit als die Links- oder
        Rechts-Endpunkt-Regel?
  \item Wie verbessern zusammengesetzte Methoden die Genauigkeit, und was ist der
        Trade-off?
  \item Warum versagen deterministische Quadraturmethoden in hohen Dimensionen?
  \item Inwiefern ist die Mittelpunktsregel ein Spezialfall der Monte-Carlo-Schätzung?
  \item Was ist die Intuition dahinter, dass Monte Carlo in hohen Dimensionen nicht
        ``explodiert''?
  \item Wann würden Sie die Simpson-Regel gegenüber Monte Carlo verwenden?
  \item Wie hängt der Mini-Batch-Mittelwert in SGD mit der Monte-Carlo-Schätzung
        zusammen?
  \item Wie hilft die Gradientenschätzung in SGD, lokalen Minima zu entkommen?
\end{itemize}

\teaser{Warum werden hochgradige Quadraturmethoden instabil, wenn sie Polynome
niedrigen Grades approximieren?}

% =====================================================================
\kapitel{Kapitel 7 \;--\; Modelltraining (Training a Model from First Principles)}

\aufg{7.1 \;--\; Loss bei Initialisierung}{Handrechnung}
Lineare Regression $\hat y(x)=wx+b$, MSE-Loss
$L(w,b)=\tfrac1n\sum_i(wx_i+b-y_i)^2$, Initialisierung $w=0$, $b=\bar y$.
Zeigen Sie, dass $L(0,\bar y)=\mathrm{Var}(y)$. Mit $\bar y=231/14=16{,}5$ und
$\sum_i(y_i-\bar y)^2=1401{,}5$: berechnen Sie den erwarteten ersten Loss-Wert. Was
bedeutet ein gedruckter Wert von $5000$ bzw.\ $0{,}0001$?

\aufg{7.2 \;--\; Overfitting als Sanity-Check}{Konzept}
Ein Grad-4-Polynom wird durch zwei Trainingspunkte (Italien $(4{,}0;4)$, UK
$(9{,}5;21)$) gelegt. Wie viele Null-Loss-Kurven gehen durch die zwei Punkte, und was
sagt diese Anzahl über die Geometrie der Verlustlandschaft aus?

\aufg{7.3 \;--\; Mini-Batch-SGD}{Konzept}
In welchem präzisen Sinn ist Mini-Batch-SGD ein Monte-Carlo-Schätzer (was wird worüber
gemittelt)? Ab welchem $n$ kippt der Engpass von voll-batch zu mini-batch, gegeben der
Fehler skaliert mit $1/\sqrt{B}$?

\aufg{7.4 \;--\; Mehrere Random Seeds}{Konzept}
Für das Modell $\hat y(x)=w^2x+b$: Warum ist der Ursprung ein Sattel statt eines
Minimums, und warum ist die Null-Initialisierung trotzdem eine Falle? Wozu dient das
Training aus mehreren Random Seeds?

\aufg{7.5 \;--\; Quantisierung für Inferenz}{Konzept}
Auf dem Schokoladen-Datensatz erreichte die Integer-Präzision fast denselben Fehler wie
volle Präzision. Warum versenden Produktionsmodelle dennoch in float32? Warum kann ein
in hoher Präzision trainiertes Modell in niedriger Präzision eingesetzt werden, während
Training in niedriger Präzision scheitert?

\reflexion
\begin{itemize}
  \item Warum entdeckt das Treiben des Loss auf $0$ bei einem überparametrisierten
        Modell Bugs in der Pipeline statt im Optimierer?
  \item Was sagt ein gedruckter erster Loss von $5000$ bzw.\ $0{,}0001$ (statt
        $\approx100$) aus?
  \item In welchem präzisen Sinn ist Mini-Batch-SGD ein Monte-Carlo-Schätzer?
  \item Warum ist für $\hat y=w^2x+b$ der Ursprung ein Sattel, und warum ist
        Null-Init eine praktische Falle?
  \item Warum lässt sich ein hochpräzise trainiertes Modell niedrigpräzise einsetzen,
        während Training in niedriger Präzision scheitert?
\end{itemize}

\end{document}
